% Worked example, version B: a storyboard, one row per step: in words | what it looks like | as a formula.
\documentclass[border=3pt]{standalone}
% Shared preamble of the worked-example figures (case_a/b/c.tex); data from ../generated/case_data.tex
\usepackage{libertinus}
\renewcommand{\familydefault}{\sfdefault}
\usepackage{tikz}
\usepackage{amsmath}
\usepackage{booktabs}
\usetikzlibrary{arrows.meta,positioning,calc,fit,backgrounds,decorations.pathreplacing}
\makeatletter
\newcommand{\C}[1]{\@ifundefined{case@#1}{\textcolor{red}{[#1?]}}{\csname case@#1\endcsname}}
\makeatother
% word-class palette (decision_shares.png / playground), one slot per category
\definecolor{catF}{HTML}{2A78D6}  % function words
\definecolor{catV}{HTML}{EB6834}  % verbs
\definecolor{catN}{HTML}{1BAF7A}  % nouns and realia
\definecolor{catA}{HTML}{EDA100}  % agents and deities
\definecolor{catX}{HTML}{E87BA4}  % foreign words
\definecolor{catO}{HTML}{9A9A94}  % other / unclassified
\definecolor{shade}{HTML}{EB6834}
\definecolor{ink}{HTML}{2B2B2B}
\definecolor{soft}{HTML}{8A8A8A}
\definecolor{topc}{HTML}{1F7A6B}
\definecolor{secc}{HTML}{5B4B8A}
% a word tinted by its class; dashed frame when the class is low-confidence
\tikzset{wc-low/.style={densely dashed}, wc-ok/.style={}, reset/.style={text width={}, minimum height=0pt, minimum width=0pt, align=center, draw=none, fill=none, solid, font=}}
\newcommand{\wc}[3]{\tikz[baseline=(w.base)]\node[reset, inner sep=0.8pt, rounded corners=1pt, fill=#1!22,
  draw=#1!70, line width=0.4pt, wc-#2](w){\strut #3};}
% a word shaded by its share of the decision (0-100)
\newcommand{\ws}[2]{\tikz[baseline=(w.base)]\node[reset, inner sep=0.8pt, rounded corners=1pt, fill=shade!#1](w){\strut #2};}
\newcommand{\broken}[1]{\textcolor{soft}{#1}}
\newcommand{\tline}[2]{\noindent\hangindent=5.7mm\hangafter=1\makebox[4.5mm][r]{\textcolor{soft}{#1}}\hspace{1.2mm}#2\par}
\newcommand{\wf}[1]{\textbf{#1}}          % a word form
\newcommand{\sg}[1]{\textit{#1}}          % a sign sequence
\newcommand{\stepno}[1]{\tikz[baseline=-0.6ex]\node[reset, circle, fill=ink, text=white, inner sep=0.6pt, font=\scriptsize\bfseries]{#1};}
\newcommand{\legend}{%
  \wc{catF}{ok}{function} \wc{catV}{ok}{verb} \wc{catN}{ok}{noun, realia} \wc{catA}{ok}{agent, deity}
  \wc{catX}{ok}{foreign} \wc{catO}{low}{unclear}}
% generated by make_case_figures.py; do not edit
\expandafter\def\csname case@txtid\endcsname{KBo 41.22}
\expandafter\def\csname case@filed\endcsname{470}
\expandafter\def\csname case@words\endcsname{33}
\expandafter\def\csname case@top\endcsname{Luwian rituals (Puriyanni, Kuwattalla)}
\expandafter\def\csname case@topshort\endcsname{Luwian rituals}
\expandafter\def\csname case@second\endcsname{Kizzuwatna rituals}
\expandafter\def\csname case@secondshort\endcsname{Kizzuwatna rituals}
\expandafter\def\csname case@s1\endcsname{0.39}
\expandafter\def\csname case@s2\endcsname{0.25}
\expandafter\def\csname case@margin\endcsname{0.14}
\expandafter\def\csname case@tau\endcsname{0.018}
\expandafter\def\csname case@nfeat\endcsname{177}
\expandafter\def\csname case@dim\endcsname{25{,}372}
\expandafter\def\csname case@ntop\endcsname{10}
\expandafter\def\csname case@nsecond\endcsname{165}
\expandafter\def\csname case@common\endcsname{−21\%}
\expandafter\def\csname case@idf.rare\endcsname{6.8}
\expandafter\def\csname case@idf.common\endcsname{1.2}
\expandafter\def\csname case@nn1\endcsname{KBo 9.125+}
\expandafter\def\csname case@nn1cth\endcsname{409}
\expandafter\def\csname case@nn.409\endcsname{5}
\def\CaseText{\tline{3′}{\wc{catJ}{ok}{i-da-a-lu-un} \wc{catO}{low}{UD-an} \wc{catO}{low}{ma-}}\tline{4′}{\wc{catA}{ok}{DINGIR\textsuperscript{MEŠ}-aš} \wc{catT}{ok}{kar-pí-in} \wc{catO}{low}{pa-an-ga-u-wa-aš} \wc{catO}{low}{EME-an}}\tline{5′}{\wc{catV}{ok}{ša-me-ed-du-wa-ra-at-kán} \broken{ḫi-x}}\tline{6′}{\wc{catF}{ok}{na-aš-ta} \wc{catA}{ok}{EN} \wc{catN}{ok}{SÍSKUR} \wc{catT}{ok}{ŠUM-ŠU}}\tline{7′}{\wc{catF}{ok}{le-e} \wc{catO}{low}{ku-wa-pí-ik-ki} \broken{a-x}}\tline{8′}{\wc{catT}{ok}{te-eš-ḫa-an} \wc{catO}{low}{ma-ni-ku-wa-an(-)}}\tline{9′}{\wc{catF}{ok}{nu-uš-ši} \wc{catN}{ok}{\textsuperscript{GI}pát-tar} \wc{catF}{ok}{pa-ra-a}}\tline{10′}{\wc{catA}{ok}{\textsuperscript{MUNUS}ŠU.GI-ma} \wc{catV}{ok}{me-ma-i} \broken{x}}\tline{11′}{\wc{catA}{ok}{DINGIR\textsuperscript{MEŠ}-aš} \wc{catT}{ok}{kar-pí-in} \wc{catO}{low}{pa-an-ga-u-wa-aš}}\tline{12′}{\wc{catA}{ok}{\textsuperscript{LÚ}A.ZU-ma-aš-ši-iš-ša-an}}\tline{13′}{\wc{catF}{ok}{na-aš} \wc{catF}{ok}{ša-ra-a} \wc{catO}{low}{ti-}}\tline{14′}{\wc{catF}{ok}{na-aš-ta} \wc{catB}{ok}{GU₄-un}}\tline{15′}{\wc{catF}{ok}{a-ra-aḫ-za-an-da}}}
\def\CaseShaded{\tline{3′}{\ws{27}{i-da-a-lu-un} \ws{2}{UD-an} \ws{0}{ma-}}\tline{4′}{\ws{2}{DINGIR\textsuperscript{MEŠ}-aš} \ws{39}{kar-pí-in} \ws{70}{pa-an-ga-u-wa-aš} \ws{93}{EME-an}}\tline{5′}{\ws{0}{ša-me-ed-du-wa-ra-at-kán} \broken{ḫi-x}}\tline{6′}{\ws{0}{na-aš-ta} \ws{0}{EN} \ws{0}{SÍSKUR} \ws{14}{ŠUM-ŠU}}\tline{7′}{\ws{0}{le-e} \ws{15}{ku-wa-pí-ik-ki} \broken{a-x}}\tline{8′}{\ws{0}{te-eš-ḫa-an} \ws{6}{ma-ni-ku-wa-an(-)}}\tline{9′}{\ws{0}{nu-uš-ši} \ws{0}{\textsuperscript{GI}pát-tar} \ws{2}{pa-ra-a}}\tline{10′}{\ws{27}{\textsuperscript{MUNUS}ŠU.GI-ma} \ws{0}{me-ma-i} \broken{x}}\tline{11′}{\ws{2}{DINGIR\textsuperscript{MEŠ}-aš} \ws{39}{kar-pí-in} \ws{70}{pa-an-ga-u-wa-aš}}\tline{12′}{\ws{0}{\textsuperscript{LÚ}A.ZU-ma-aš-ši-iš-ša-an}}\tline{13′}{\ws{0}{na-aš} \ws{0}{ša-ra-a} \ws{5}{ti-}}\tline{14′}{\ws{0}{na-aš-ta} \ws{0}{GU₄-un}}\tline{15′}{\ws{0}{a-ra-aḫ-za-an-da}}}
\def\CaseKeyWord{\wc{catO}{low}{pa-an-ga-u-wa-aš}}
\def\CaseBars{\simbar{0}{Luwian rituals}{0.389}{top}\simbar{1}{Kizzuwatna rituals}{0.245}{second}\simbar{2}{Hattian cult layer}{0.204}{other}\simbar{3}{Syro-Mesopotamia}{0.181}{other}\simbar{4}{Kizzuwatna–Lower Land}{0.173}{other}\simbar{5}{Palaic festivals}{0.165}{other}\simbar{6}{Luwian festivals}{0.165}{other}\simbar{7}{Ḫuwaššanna festivals}{0.160}{other}\simbar{8}{Mixed Hurrian}{0.147}{other}\simbar{9}{Arzawa}{0.146}{other}}
\def\CaseFeatRows{\wf{EMEan} & 0.11 & 0.17 & 0.00 & +18.1 \\\sg{EME-an} & 0.10 & 0.17 & 0.00 & +17.1 \\\wf{pankauaš} & 0.18 & 0.08 & 0.00 & +14.7 \\\sg{ga-u-wa} & 0.17 & 0.08 & 0.00 & +12.7 \\\sg{an-ga-u} & 0.16 & 0.07 & 0.00 & +11.8 \\\sg{pa-an-ga} & 0.16 & 0.07 & 0.00 & +11.8 \\\midrule \wf{EN} & 0.05 & 0.00 & 0.06 & −3.0 \\\sg{EN} & 0.04 & 0.00 & 0.06 & −2.3 \\}
\def\CaseTug{\tug{0}{\wf{EMEan}}{18.1}\tug{1}{\sg{EME-an}}{17.1}\tug{2}{\wf{pankauaš}}{14.7}\tug{3}{\sg{ga-u-wa}}{12.7}\tug{4}{\sg{an-ga-u}}{11.8}\tug{5}{\sg{pa-an-ga}}{11.8}\tug{6}{\wf{EN}}{-3.0}\tug{7}{\sg{EN}}{-2.3}}
\expandafter\def\csname case@x.w\endcsname{0.18}
\expandafter\def\csname case@x.n1\endcsname{0.16}
\expandafter\def\csname case@x.n2\endcsname{0.16}
\expandafter\def\csname case@x.n3\endcsname{0.17}


\begin{document}
\begin{tikzpicture}[font=\sffamily\scriptsize\color{ink},
  words/.style={text width=34mm, anchor=north west, inner sep=2pt, align=left},
  visual/.style={text width=78mm, anchor=north west, inner sep=2pt},
  formula/.style={text width=36mm, anchor=north west, inner sep=3pt, fill=soft!8, rounded corners=2pt},
  head/.style={font=\sffamily\scriptsize\bfseries\color{soft}, anchor=south west, inner sep=1pt}]
\def\xw{0mm}\def\xv{37mm}\def\xf{118mm}
\node[head] at (\xw,1mm) {IN WORDS}; \node[head] at (\xv,1mm) {ON THE FRAGMENT}; \node[head] at (\xf,1mm) {AS A FORMULA};
% row 1
\node[words] (w1) at (\xw,0) {\stepno{1} \textbf{Read.} An orphan, \C{txtid}, filed under CTH \C{filed}. Each legible word gets a word class from the edition's analysis; dashed: unclear.};
\node[visual] (v1) at (\xv,0) {{\linespread{1.15}\selectfont\raggedright\CaseText}\vspace{1pt}{\tiny\legend}};
\node[formula] (f1) at (\xf,0) {\C{words} legible words.};
\path let \p1=(w1.south), \p2=(v1.south), \p3=(f1.south) in coordinate (r1) at (0,{min(\y1,\y2,\y3)-2mm});
\draw[soft!50] (\xw, 0 |- r1) -- (\xf+38mm, 0 |- r1);
% row 2
\node[words] (w2) at (\xw, 0 |- r1) {\stepno{2} \textbf{Count.} Every word form and every sign sequence inside a word is a feature. A feature counts more when it is frequent in the fragment and rare in the corpus.};
\node[visual] (v2) at (\xv, 0 |- r1) {%
  \wc{catN}{ok}{pa-an-ga-u-wa-aš} $\to$ \wf{pankauaš} \C{x.w}, \sg{pa-an-ga} \C{x.n1}, \sg{an-ga-u} \C{x.n2}, \sg{ga-u-wa} \C{x.n3}\\[3pt]
  Rare word \wf{pankauaš}: idf \C{idf.rare}. Common sign \sg{na}: idf \C{idf.common}.\\[2pt]
  The fragment becomes a vector $\mathbf{x}$ with \C{nfeat} non-zero entries out of \C{dim}.};
\node[formula] (f2) at (\xf, 0 |- r1) {$x_f \propto (1+\log \mathit{tf}_f)\,\mathit{idf}_f$\\[2pt] $\mathit{idf}_f=\log\frac{1+N}{1+\mathit{df}_f}+1$\\[2pt] $\lVert\mathbf{x}\rVert = 1$};
\path let \p1=(w2.south), \p2=(v2.south), \p3=(f2.south) in coordinate (r2) at (0,{min(\y1,\y2,\y3)-2mm});
\draw[soft!50] (\xw, 0 |- r2) -- (\xf+38mm, 0 |- r2);
% row 3
\node[words] (w3) at (\xw, 0 |- r2) {\stepno{3} \textbf{Compare.} Each tradition is the average of its manuscripts. The fragment is compared with each average.};
\node[visual] (v3) at (\xv, 0 |- r2) {%
  \begin{tikzpicture}[x=55mm, y=2.9mm]
    \newcommand{\simbar}[4]{%
      \def\k{#4}\def\kt{top}\def\ks{second}%
      \ifx\k\kt\def\col{topc}\else\ifx\k\ks\def\col{secc}\else\def\col{soft!45}\fi\fi
      \fill[\col] (0,-#1) rectangle (#3,-#1-0.72);
      \node[anchor=east, inner sep=1pt, font=\tiny] at (0,-#1-0.36) {#2};
      \node[anchor=west, inner sep=1pt, font=\tiny] at (#3,-#1-0.36) {#3};}
    \CaseBars
  \end{tikzpicture}};
\node[formula] (f3) at (\xf, 0 |- r2) {$\mathbf{c}_t=\dfrac{\sum_{d\in t}\mathbf{x}_d}{\lVert\sum_{d\in t}\mathbf{x}_d\rVert}$\\[4pt] $s_t=\cos(\mathbf{x},\mathbf{c}_t)=\sum_f x_f\,c_{t,f}$};
\path let \p1=(w3.south), \p2=(v3.south), \p3=(f3.south) in coordinate (r3) at (0,{min(\y1,\y2,\y3)-2mm});
\draw[soft!50] (\xw, 0 |- r3) -- (\xf+38mm, 0 |- r3);
% row 4
\node[words] (w4) at (\xw, 0 |- r3) {\stepno{4} \textbf{Decide.} Place the fragment only if the nearest tradition leads the runner-up clearly; otherwise leave it unassigned.};
\node[visual] (v4) at (\xv, 0 |- r3) {%
  \begin{tikzpicture}[x=55mm, y=3mm]
    \fill[topc] (0,0) rectangle (\C{s1},0.8); \node[anchor=east, font=\tiny] at (0,0.4) {\C{topshort}};
    \fill[secc] (0,-1.2) rectangle (\C{s2},-0.4); \node[anchor=east, font=\tiny] at (0,-0.8) {\C{secondshort}};
    \draw[decorate, decoration={brace, amplitude=3pt}, line width=0.6pt] (\C{s1},1.1) -- (\C{s2},1.1) node[pos=0.5, above=3pt, font=\tiny] {};
    \draw[densely dotted] (\C{s2},-1.2) -- (\C{s2},1.1);
    \node[anchor=west, font=\scriptsize] at (\C{s1}+0.02,0) {$m=\C{margin} \geq \C{tau}$ $\Rightarrow$ \textbf{\color{topc}\C{topshort}}};
  \end{tikzpicture}};
\node[formula] (f4) at (\xf, 0 |- r3) {$m=s_1-s_2=\C{s1}-\C{s2}$\\[2pt] place if $m\geq\tau=\C{tau}$ \\ ($\tau$: 80\% of test fragments right)};
\path let \p1=(w4.south), \p2=(v4.south), \p3=(f4.south) in coordinate (r4) at (0,{min(\y1,\y2,\y3)-2mm});
\draw[soft!50] (\xw, 0 |- r4) -- (\xf+38mm, 0 |- r4);
% row 5
\node[words] (w5) at (\xw, 0 |- r4) {\stepno{5} \textbf{Explain.} The lead is a sum over features, so each word's share of the decision can be shown. Here the formula \sg{pa-an-ga-u-wa-aš EME-an} decides.};
\node[visual] (v5) at (\xv, 0 |- r4) {{\linespread{1.15}\selectfont\raggedright\CaseShaded}\vspace{1pt}{\tiny Darker: larger share of the lead.}};
\node[formula] (f5) at (\xf, 0 |- r4) {$m=\sum_f \underbrace{x_f\,(c_{1,f}-c_{2,f})}_{\text{share of feature }f}$};
\path let \p1=(w5.south), \p2=(v5.south), \p3=(f5.south) in coordinate (r5) at (0,{min(\y1,\y2,\y3)-2mm});
\draw[soft!50] (\xw, 0 |- r5) -- (\xf+38mm, 0 |- r5);
% row 6
\node[words] (w6) at (\xw, 0 |- r5) {\stepno{6} \textbf{Check by hand.} The closest single manuscripts are candidates for the philologist.};
\node[visual] (v6) at (\xv, 0 |- r5) {Of the five nearest manuscripts, \C{nn.409} are Tunnawiya's ritual (CTH \C{nn1cth}), which no reference group includes.};
\node[formula] (f6) at (\xf, 0 |- r5) {$\arg\max_d \cos(\mathbf{x},\mathbf{x}_d)$};
\end{tikzpicture}
\end{document}
